\begin{definition}
Fix the graph, palette, distinguished edge, exact terminal tree \(T\),
generating coloring \(\varphi_n\), positive-rung ETT, hierarchy base,
minimum hierarchy counterexample \(C\), and its
\noderef{HierarchyMinimumCounterexampleContext}.  A \emph{common
Situation 7.2 tail datum} consists of the following exact data, before the
split according to whether the last-edge color is missing at the penultimate
vertex.  Put \(\ell=\noderef{HierarchyFullTailLength}(C)\).  One has
\[
 \ell\geq2,
 \qquad
 \noderef{HierarchyPathNumber}(C.\mathit{hierarchy})=\ell .
\]
There are indices \(s\in\operatorname{Fin}|T|\) and
\(a,b\in\operatorname{Fin}(|T|+1)\) satisfying
\[
 s+1=|T|,\qquad a+1=|T|,\qquad b=|T|.
\]
Writing \(e_p=T.\mathit{edge}(s)\),
\(y_{p-1}=T.\mathit{vertex}(a)\), and
\(y_p=T.\mathit{vertex}(b)\), the individual edge \(e_p\) is not incident
with \(y_{p-1}\).  Finally a color \(\theta\) is selected with
\(\varphi_n(e_p)=\theta\).  No assertion is made here about whether
\(\theta\) is missing at \(y_{p-1}\); that is precisely the later split
between Subcases 3.1 and 3.2.
\end{definition}
